%---------------------------------------------------------------------
% Title page and one orientation page.
%---------------------------------------------------------------------
\thispagestyle{empty}

% Band depth is measured from the rendered page, never computed: Nastaliq at
% Scale=1.45 sets much taller lines than the nominal point sizes suggest, and
% every masthead here has run deeper than a first estimate. Leave headroom --
% empty band below the last line looks fine; a clipped line does not.
\begin{tikzpicture}[remember picture, overlay]
  \fill[bmaccent] ([yshift=-12.3cm]current page.north west) rectangle
        ([yshift=-12.5cm]current page.north east);
\end{tikzpicture}

\vspace*{0.9cm}
{\color{bmink}
  {\large علامہ اقبال اوپن یونیورسٹی، اسلام آباد}\par
  \vspace{3pt}
  {شعبۂ ابلاغِ عامہ \ \textbullet\ \ بی اے}\par
  \vspace{20pt}
  {\Huge\bfseries\color{bmblue} اصولِ صحافت}\par
  \vspace{6pt}
  {\large\bfseries \textenglish{Principles of Journalism}}\par
  \vspace{12pt}
  {\LARGE\bfseries متوقع سوالات اور ان کے جوابات}\par
  \vspace{12pt}
  {\large کورس کوڈ \textenglish{430}}\par
  \vspace{4pt}
  % \par inside the group: colour applies when the paragraph breaks.
  {چودہ سوالات \ \textbullet\ \ دو حقیقی پرچوں سے ماخوذ\par}
}

\vspace{1.0cm}

\begin{tcolorbox}[enhanced, colback=white, colframe=bmblue, boxrule=0.6pt,
                  arc=3pt, left=12pt, right=12pt, top=10pt, bottom=10pt]
\textbf{\color{bmblue}سوالات کہاں سے لیے گئے ہیں؟}\ \
اس کتابچے کے تمام سوالات اے آئی او یو کے \textbf{دو حقیقی گزشتہ پرچوں} سے
لفظ بہ لفظ نقل کیے گئے ہیں:

\smallskip
\begin{itemize}\setlength{\itemsep}{0pt}\setlength{\topsep}{2pt}
  \item بی اے، اصولِ صحافت (\textenglish{430})، سمسٹر \textbf{بہار
        \textenglish{2013}} --- آٹھ سوالات
  \item بی اے، اصولِ صحافت (\textenglish{430})، سمسٹر \textbf{خزاں
        \textenglish{2013}} --- سات سوالات
\end{itemize}

\smallskip
دونوں پرچوں میں دہرائے جانے والے سوالات کو یکجا کر دیا گیا ہے --- اور یہی
دہرایا جانا بتاتا ہے کہ کون سے موضوعات سب سے اہم ہیں۔

\medskip
\textbf{\color{bmblue}ایک صاف بات۔}\ \
اس کورس کی سرکاری درسی کتاب دستیاب نہیں ہو سکی، اس لیے یہ کتابچہ کتاب کے
اختتامی سوالات پر نہیں بلکہ \emph{گزشتہ پرچوں} پر مبنی ہے۔ جوابات صحافت کے
مسلمہ اصولوں اور پاکستانی صحافت کی معلوم تاریخ کے مطابق لکھے گئے ہیں۔ جہاں
کتاب کی مخصوص تفصیل درکار ہو، وہاں اپنی درسی کتاب سے ملا لیجیے۔

\medskip
\textbf{\color{bmblue}حیثیت۔}\ \
ذاتی مطالعے کے لیے۔ یہ اے آئی او یو کی مطبوعہ کتاب نہیں۔ اسے من و عن نقل کر کے
اسائنمنٹ میں جمع کرانا منع ہے۔
\end{tcolorbox}

\clearpage

%---------------------------------------------------------------------
\section*{پرچے کا خاکہ اور طریقۂ تیاری}
\markboth{پرچے کا خاکہ}{}

دونوں پرچوں سے پرچے کی ساخت بالکل واضح ہے، اور یہ جاننا آدھی تیاری ہے:

\medskip
\begin{center}
\begin{tabular}{@{}p{4.6cm} p{9.8cm}@{}}
\toprule
کل نمبر & \textenglish{100} \\
کامیابی کے نمبر & \textenglish{40} \\
وقت & تین گھنٹے \\
سوالات & سات یا آٹھ دیے جاتے ہیں \\
کتنے حل کرنے ہیں & \textbf{کوئی سے پانچ} \\
فی سوال نمبر & \textenglish{20} --- تمام سوالات کے نمبر مساوی ہیں \\
\bottomrule
\end{tabular}
\end{center}

\medskip
\begin{nukta}[بیس نمبر کا جواب کیسا ہونا چاہیے]
تین گھنٹوں میں پانچ سوال، یعنی فی سوال تقریباً \textbf{پینتیس منٹ}۔ ہر جواب کا
سانچہ یہ رکھیے:

\smallskip
\textbf{(۱) تعریف} --- ایک یا دو سطر میں، جامع۔ \quad
\textbf{(۲) اہمیت یا پس منظر} --- ایک مختصر پیرا۔ \quad
\textbf{(۳) اصل مواد} --- اقسام، خصوصیات، فرائض یا نکات، \emph{عنوانات کے
ساتھ}۔ \quad
\textbf{(۴) پاکستانی مثالیں} --- اخبارات، شخصیات، واقعات۔ \quad
\textbf{(۵) خلاصہ} --- دو سطریں۔

\smallskip
سب سے زیادہ نمبر \emph{عنوانات} اور \emph{مثالوں} سے ملتے ہیں۔ مسلسل پیراگراف
میں لکھا گیا جواب، خواہ درست ہو، کم نمبر لیتا ہے۔
\end{nukta}

\medskip
\noindent
\textbf{\color{bmblue}کون سے موضوعات سب سے زیادہ دہرائے گئے}

\medskip
\begin{center}
\begin{tabular}{@{}p{6.4cm} p{2.6cm} p{5.4cm}@{}}
\toprule
\textbf{موضوع} & \textbf{کتنے پرچوں میں} & \textbf{اس کتابچے میں سوال} \\
\midrule
صحافت کا ارتقاء (برصغیر و پاکستان) & دونوں & \textenglish{2} \\
سب ایڈیٹر کے اوصاف و فرائض        & دونوں & \textenglish{7} \\
ضابطۂ اخلاق اور صحافتی قوانین      & دونوں & \textenglish{8}، \textenglish{9} \\
اداریہ --- تعریف، اقسام، فرق       & دونوں & \textenglish{4}، \textenglish{5} \\
صحافت کی تعریف اور نظریات          & بہار  & \textenglish{1} \\
کالم نویسی                          & بہار  & \textenglish{6} \\
آزادیٔ صحافت                        & بہار  & \textenglish{10} \\
سر سید اور ظفر علی خان              & خزاں  & \textenglish{3} \\
جدید رجحانات                        & خزاں  & \textenglish{12} \\
صحافت اور جمہوریت                   & خزاں  & \textenglish{11} \\
\bottomrule
\end{tabular}
\end{center}

\medskip
\noindent
پہلی چار سطریں دونوں پرچوں میں آئی ہیں۔ اگر وقت بہت کم ہو تو انہی چار سے شروع
کیجیے --- پانچ میں سے تین سوال عموماً انہی سے بن جاتے ہیں۔
